\documentclass[conference]{IEEEtran}
\usepackage[T2A]{fontenc}
\usepackage[utf8]{inputenc}
\usepackage[russian]{babel}
\usepackage{cite}
\usepackage{amsmath,amssymb,amsfonts}
\usepackage{algorithmic}
\usepackage{graphicx}
\usepackage{textcomp}
\usepackage{xcolor}
\usepackage{hyperref}

% Правильные переносы (добавлены для русского)
\hyphenation{ге-не-ра-тив-ный ре-гу-ля-тор-ный гор-мо-наль-ный эн-до-крин-ный}

\begin{document}

\renewcommand{\abstractname}{Аннотация}
\renewcommand{\IEEEkeywordsname}{Ключевые слова}

\title{GRA-Гормональный Взрыв Субъекта: Генеративный Регуляторный Автомат для Моделирования Эндокринных Всплесков}

\author{\IEEEauthorblockN{qqewq}
\IEEEauthorblockA{Независимый исследователь\\
GitHub: qqewq\\
\href{https://github.com/qqewq/GRA-Hormonal-Explosion-Of-Subject}{Репозиторий проекта}}}

\maketitle

\begin{abstract}
Биологические системы демонстрируют сложные гормональные всплески, управляемые обратными связями, которые регулируют критические физиологические и поведенческие переходы. В данной статье представлена модель \textbf{GRA-Hormonal Explosion of Subject} (GRA-HES) — генеративный регуляторный автомат, предназначенный для моделирования, визуализации и анализа острых эндокринных выбросов у виртуального субъекта. Модель сочетает генетический регуляторный алгоритм (GRA) с многослойным графом взаимодействий гормон–рецептор, что порождает каскадные, самоусиливающиеся взрывные события. Описаны математическая формулировка, реализация и эмерджентные паттерны, наблюдаемые при различных режимах параметров. Публикация исходного кода предоставляет исследователям воспроизводимый испытательный стенд для изучения гормезисной зависимости «доза–эффект», пороговых колебательных явлений и сценариев синтетического эндокринного разрушения.
\end{abstract}

\begin{IEEEkeywords}
генеративный регуляторный автомат, гормональный взрыв, эндокринная динамика, вычислительная биология, агентное моделирование
\end{IEEEkeywords}

\section{Введение}
Быстрые гормональные флуктуации, называемые «гормональными взрывами», лежат в основе полового созревания, стрессовых реакций и метаболических переключений \cite{m1}. Хотя модели на основе дифференциальных уравнений хорошо описывают плавные эндокринные ритмы, они часто не способны воспроизвести резкие, переключательные выбросы, наблюдаемые \textit{in vivo}. Недавние достижения в области генеративных регуляторных сетей \cite{m2} предлагают более выразительный формализм, однако их применение к эндокринным системам остаётся неисследованным.

Проект \textbf{GRA-Hormonal Explosion of Subject} заполняет этот пробел, создавая цифрового двойника субъекта, чьи внутренние железистые слои управляются Генеративным Регуляторным Автоматом (GRA). GRA эволюционирует вектор гормонов через дискретно-временные, основанные на правилах взаимодействия, обеспечивая как стохастическую базальную секрецию, так и детерминированные взрывы обратной связи. Сопутствующий визуализатор отображает состояние субъекта в виде динамического аватара, делая гормональные кризы интуитивно наблюдаемыми.

\section{Описание модели}
\subsection{Генеративный Регуляторный Автомат (GRA)}
GRA — это кортеж $\mathcal{G} = (\mathbf{H}, \mathbf{R}, \mathbf{W}, \tau)$, где:
\begin{itemize}
    \item $\mathbf{H} \in \mathbb{R}^n$ — вектор концентраций гормонов (например, кортизол, адреналин, тестостерон, эстроген).
    \item $\mathbf{R} \in \mathbb{R}^{m}$ — чувствительности рецепторов.
    \item $\mathbf{W} \in \mathbb{R}^{n \times (n+m)}$ — весовая матрица, кодирующая стимулирующие/ингибирующие взаимодействия.
    \item $\tau$ — пороговая функция, запускающая выполнение дискретных правил.
\end{itemize}
На каждом временно́м шаге $t$ состояние обновляется согласно
\begin{equation}
    \mathbf{H}_{t+1} = \mathbf{H}_t + \Delta t \cdot \big( \mathbf{W}_{H} \cdot \sigma(\mathbf{H}_t \oplus \mathbf{R}_t) + \boldsymbol{\xi}_t \big),
    \label{eq:update}
\end{equation}
где $\sigma(\cdot)$ — насыщающая нелинейность (например, $\tanh$), $\oplus$ обозначает конкатенацию, а $\boldsymbol{\xi}_t \sim \mathcal{N}(0,\Sigma)$ — гауссовский шум. \textbf{Взрыв} определяется как переходное событие, при котором $\|\mathbf{H}_t\|$ превышает заданный порог всплеска $B$ в течение не менее $k$ последовательных шагов.

\subsection{Слой субъекта и триггер гормонального взрыва}
Субъект представлен в виде конечного автомата, состояниями которого являются эмоциональные/поведенческие состояния (например, спокойствие, возбуждение, взрыв). Переходы зависят от доминирующего гормона:
\begin{equation}
    S_{t+1} = f(S_t, \operatorname{argmax}(\mathbf{H}_t)).
\end{equation}
Внешний стрессор или циркадный триггер может ввести импульс $\mathbf{p}$ в определённый гормональный канал, инициируя каскад. Если возникающая положительная обратная связь (через $\mathbf{W}$) выводит систему выше $B$, происходит \emph{гормональный взрыв субъекта} — аватар визуально «взрывается», а внутренние уровни гормонов осциллируют хаотически, прежде чем вернуться к базовому уровню.

\section{Реализация}
Программный код (\texttt{GRA-Hormonal-Explosion-Of-Subject}) написан на Python с визуализацией на Pygame. Основные модули:
\begin{itemize}
    \item \texttt{gra\_core.py}: реализует уравнение (\ref{eq:update}) с настраиваемыми сетями желёз.
    \item \texttt{subject\_layer.py}: управляет конечным автоматом субъекта и триггерами анимации.
    \item \texttt{explosion\_detector.py}: обнаруживает взрывные события и регистрирует метрики (пиковая амплитуда, длительность, время восстановления).
    \item \texttt{config.yaml}: файл параметров для $\mathbf{W}$, порогов и шума.
\end{itemize}
Симуляция выполняется в реальном времени; пользователь может возмущать субъекта, кликая для инъекции синтетических гормонов или планируя стрессоры.

\section{Результаты и эмерджентные явления}
\subsection{Сигнатуры взрывов}
При стандартных параметрах система демонстрирует три режима:
\begin{enumerate}
    \item \textbf{Спокойный}: низкоамплитудное блуждание около гомеостатической уставки.
    \item \textbf{Пульсирующий}: периодические, предсказуемые всплески, напоминающие ультрадианные ритмы.
    \item \textbf{Хаотический взрыв}: апериодические, высокоамплитудные выбросы с длинными хвостами восстановления, запускаемые пересечением критического многообразия в $\mathbf{W}$.
\end{enumerate}
На рис.~\ref{fig:timeseries} показано типичное взрывное событие. Положительная обратная связь между кортизолом и адреналином создаёт быстрый пик, за которым следует ингибирующая компенсация, приводящая к временному дефициту.

\begin{figure}[htbp]
\centerline{\includegraphics[width=\linewidth]{timeseries.png}}
\caption{Концентрации гормонов во время индуцированного взрыва. Кортизол (сплошная линия) и адреналин (пунктирная линия) демонстрируют совместное всплесковое поведение. Серая полоса обозначает интервал взрыва.}
\label{fig:timeseries}
\end{figure}

\subsection{Чувствительность к топологии сети}
Мы протестировали случайно сгенерированные матрицы $\mathbf{W}$. Сети с сильным взаимным возбуждением между двумя железами (например, мотивы гипоталамо-гипофизарно-надпочечниковой оси) были значительно более склонны к взрывам (критерий Манна–Уитни, $p < 0.001$). Более разреженные топологии прямого распространения требовали более мощных внешних импульсов для вызова события.

\section{Обсуждение}
Модель GRA-HES воспроизводит характерные черты гормональных взрывов: порогово-зависимый запуск, ответ по принципу «всё или ничего» и медленное восстановление. Её генеративная природа позволяет исследовать сценарии эндокринных нарушений «что, если» без экспериментов на животных. Публичный репозиторий (\url{https://github.com/qqewq/GRA-Hormonal-Explosion-Of-Subject}) приглашает сообщество расширять библиотеки желёз, учитывать фармакокинетические задержки или связывать субъекта с многоагентными социальными симуляциями.

Ограничения включают абстрактные, безразмерные единицы гормонов и отсутствие пространственной компартментализации. В будущей работе планируется интегрировать слой переноса гормонов (диффузия между железами) и откалибровать параметры по эмпирическим данным реакции пробуждения кортизола.

\section{Заключение}
Мы представили GRA-Hormonal Explosion of Subject — гибкую вычислительную модель острых эндокринных выбросов, построенную на генеративном регуляторном автомате. Симуляция успешно порождает реалистичную динамику взрывов и служит интерактивным образовательным и исследовательским инструментом. Репозиторий является открытым и активно поддерживается.

\section*{Благодарности}
Автор благодарит сообщество открытого программного обеспечения за базовые библиотеки, сделавшие этот проект возможным.

\begin{thebibliography}{00}
\bibitem{m1}
G. Leng and M. Ludwig, ``Information processing in the hypothalamus: the importance of pulsatile signalling,'' \textit{J. Neuroendocrinol.}, vol. 20, no. 5, pp. 477–485, 2008.
\bibitem{m2}
R. Albert and H. G. Othmer, ``The topology of the regulatory interactions predicts the expression pattern of the segment polarity genes in Drosophila melanogaster,'' \textit{J. Theor. Biol.}, vol. 223, no. 1, pp. 1–18, 2003.
\bibitem{m3}
S. H. Strogatz, \textit{Nonlinear Dynamics and Chaos}. Westview Press, 2015.
\end{thebibliography}

\end{document}